\DocumentMetadata{
  pdfversion=2.0,
  pdfstandard=ua-2,
  lang=en-US,
  tagging=on,
  tagging-setup={math/setup=mathml-SE,tabsorder=structure}
}
\documentclass[11pt,letterpaper]{article}
\usepackage[margin=0.68in,includefoot]{geometry}
\usepackage{newcomputermodern}
\usepackage{amsmath,amscd,mathtools}
\usepackage{tikz,adjustbox}
\providecommand{\square}{\mdlgwhtsquare}
\usepackage[hidelinks]{hyperref}
\hypersetup{
  pdftitle={Analysis First-Year Exam, May 2018, Part 1},
  pdfauthor={University of Florida Department of Mathematics},
  pdfsubject={Graduate mathematics examination},
  pdfdisplaydoctitle=true,
  bookmarksopen=true
}
\setcounter{secnumdepth}{0}
\setlength{\parindent}{0pt}
\setlength{\parskip}{0.28em}
\setlength{\emergencystretch}{3em}
\setlength{\leftmargini}{2.15em}
\setlength{\leftmarginii}{2.65em}
\setlength{\leftmarginiii}{3em}
\renewcommand{\labelenumi}{\textbf{\theenumi.}}
\pagestyle{empty}
\raggedbottom
\allowdisplaybreaks
\newenvironment{examproblems}[1]{%
  \begin{enumerate}%
  \setcounter{enumi}{#1}%
  \setlength{\topsep}{0.35em}%
  \setlength{\partopsep}{0pt}%
  \setlength{\itemsep}{0.48em}%
  \setlength{\parsep}{0.18em}%
}{\end{enumerate}}
\newenvironment{examsubparts}{%
  \begin{enumerate}%
  \setlength{\topsep}{0.18em}%
  \setlength{\partopsep}{0pt}%
  \setlength{\itemsep}{0.16em}%
  \setlength{\parsep}{0.1em}%
}{\end{enumerate}}
\newenvironment{examinstructions}{%
  \par\begingroup\itshape
}{%
  \par\endgroup\vspace{0.18em}
}
\makeatletter
\renewcommand\section{\@startsection{section}{1}{0pt}%
  {-1.25ex plus -.25ex minus -.1ex}%
  {0.55ex plus .1ex}%
  {\normalfont\large\bfseries}}
\renewcommand{\@maketitle}{%
  \newpage\null\vskip -1em%
  {\footnotesize\bfseries
    \href{https://www.ufl.edu/}{University of Florida}, \href{https://math.ufl.edu/}{Department of Mathematics}\par}%
  \vskip -0.5em%
  \tagstructbegin{tag=Title}%
    {\LARGE\bfseries\href{https://gma.math.ufl.edu/exams/analysis-first-year/}{Analysis First-Year Exam}, May 2018, Part 1\par}%
  \tagstructend%
  \vskip 0.25em%
  \tagmcbegin{artifact}%
    \hrule height 1pt%
  \tagmcend%
  \vskip 1em%
}
\makeatother
\title{Analysis First-Year Exam, May 2018, Part 1}
\author{}
\date{}
\begin{document}
\maketitle
\thispagestyle{empty}
\begin{examinstructions}
Answer FOUR questions in detail. State carefully any results used without proof.
\end{examinstructions}
\begin{examproblems}{0}
\item
Let \(f:[0,1]\to[0,1]\) be a (weakly) increasing function. Prove that 
\[
\alpha:=\inf\{t\in[0,1]:f(t)\leq t\}
\]
 exists and is fixed by~\(f\). Show further that~\(\alpha\) is the smallest fixed point of~\(f\).
\item
Let~\(M\) be a metric space, let \(U\subseteq M\) be open and let \(A\subseteq M\) be arbitrary. Prove that \(\overline{A}\cap U\subseteq\overline{A\cap U}\). Can this inclusion be strict? Explain.
\item
Let \((a_n:n\geq 0)\) be a sequence of real numbers. Assume that there exists \(\lambda\in(0,1)\) such that \(|a_{n+1}-a_n|\leq\lambda|a_n-a_{n-1}|\) for each positive integer~\(n\). Prove that the sequence \((a_n:n\geq 0)\) converges.
\item
Let \(f:X\to Y\) be a continuous bijection between metric spaces. For each of the following statements, decide whether true or false, giving proof or counterexample as appropriate:
\begin{examsubparts}
\item[\textnormal{(i)}]
if~\(X\) is compact, then \(f^{-1}:Y\to X\) is continuous;
\item[\textnormal{(ii)}]
if~\(Y\) is compact, then \(f^{-1}:Y\to X\) is continuous.
\end{examsubparts}
\item
Let the real-valued function~\(f\) be differentiable on the open interval \((0,1)\). Show that if its derivative \(f'\) is bounded, then~\(f\) extends to a continuous function on the closed interval \([0,1]\).
\end{examproblems}
\end{document}
